% Source: Wald GR, physical PDF pp. 431--433
%         (printed pp. 418--420).
% Coverage: Chapter 14 problems 1--7.
% Figures/tables/footnotes: none.
% Status: source-reviewed and compile-clean.
% Modernization: all script symbols use \mathcal.
% Uncertainties: none.

\chapterproblems

\begin{enumerate}[label=\arabic*.,leftmargin=2em,itemsep=1.2ex]
\item
Let \(V\) be the collection of infinite sequences of complex numbers
\(\{a_i\}=(a_1,a_2,\ldots)\) such that only \emph{finitely} many of the \(a_i\)
are nonzero. Define addition and scalar multiplication on \(V\) by
\(\{a_i\}+\{b_i\}=\{a_i+b_i\}\) and \(c\{a_i\}=\{ca_i\}\) to make \(V\) a
vector space over \(\mathbb C\). Define
\[
  (\{a_i\},\{b_i\})=\sum_{i=1}^{\infty}\overline a_i b_i.
\]
\begin{enumerate}[label=\alph*),leftmargin=2em,itemsep=.45ex]
\item Show that \((\ ,\ )\) is an inner product, thus making \(V\) an inner product
space.
\item Show that \(V\) is incomplete and thus is not a Hilbert space.
\item Let \(V'\) be the collection of sequences satisfying
\(\sum_{i=1}^{\infty}|a_i|^2<\infty\). Define a vector space structure and inner
product on \(V'\) in the same way as for \(V\). Show that \(V'\) is complete and
thus defines a Hilbert space, usually denoted \(\ell_2\). (You may use the fact that
\(\mathbb C\) is complete in your proof.)
\end{enumerate}

\item
\begin{enumerate}[label=\alph*),leftmargin=2em,itemsep=.45ex]
\item Show that the annihilation and creation operators defined by
equations~\eqref{eq:14.2.10} and~\eqref{eq:14.2.12} satisfy the commutation
relations
\begin{align*}
  [a(\overline{\sigma}),a(\rho)]&=0,\\
  [a^\dagger(\sigma),a^\dagger(\rho)]&=0,\\
  [a(\overline{\sigma}),a^\dagger(\rho)]&=(\sigma,\rho)I
\end{align*}
for all \(\sigma,\rho\in\mathcal H\), where the commutator of two operators is
defined by \([A,B]=AB-BA\), and \(I\) is the identity operator on
\(\mathcal F_s(\mathcal H)\).
\item Show that the operator
\[
  N(\sigma)=a^\dagger(\sigma)a(\overline{\sigma})
\]
can be interpreted as the number operator for the (normalized) state
\(\sigma\in\mathcal H\), i.e., the eigenstates of \(N(\sigma)\) are the states with
a definite number of particles in state \(\sigma\), and the eigenvalues of
\(N(\sigma)\) in these eigenstates are the number of such particles.
\end{enumerate}

\item
Define the operators
\(A\colon\mathcal H_{\mathrm{out}}\longrightarrow\mathcal H_{\mathrm{in}}\),
\(B\colon\mathcal H_{\mathrm{out}}\longrightarrow
\overline{\mathcal H}_{\mathrm{in}}\) by the \lq\lq time reverse\rq\rq\ of the
definitions of \(C\) and \(D\) given in the text, i.e., for
\(\tau\in\mathcal H_{\mathrm{out}}\), \(A\tau\) is the state in
\(\mathcal H_{\mathrm{in}}\) obtained by propagating into the past the solution
\(\widetilde\tau\) which agrees with \(\tau\) in the future and taking its positive
frequency part in the past, while \(B\tau\) is the state in
\(\overline{\mathcal H}_{\mathrm{in}}\) associated with its negative frequency part
in the past. Use the independence of the Klein--Gordon inner product,
equation~\eqref{eq:14.2.5}, on the choice of Cauchy surface together with the fact
that in the asymptotic future or past the Klein--Gordon inner product of any positive
frequency solution with any negative frequency solution vanishes to prove that the
following relations are satisfied by \(A,B,C,D\):
\begin{align*}
  A^\dagger A-B^\dagger B&=I,\\
  B^\dagger\overline A&=A^\dagger\overline B,\\
  C^\dagger C-D^\dagger D&=I,\\
  D^\dagger\overline C&=C^\dagger\overline D,\\
  A^\dagger&=C,\qquad B^\dagger=-\overline D.
\end{align*}
Note that the relation \(D^\dagger\overline C=C^\dagger\overline D\) implies
that \(E=\overline D\,\overline C^{-1}\) is symmetric,
\(E^\dagger=\overline E\).

\item
Consider a conformally invariant field (see appendix~D) in a spacetime of the form
illustrated in Figure~14.1 which, in addition, is conformally flat, i.e.,
\(g_{ab}=\Omega^2\eta_{ab}\) everywhere. Show that no particle creation can occur.

\item
With \(A\) and \(B\) defined as in problem~3, it is clear that by \lq\lq time
reversal symmetry\rq\rq\ equation~\eqref{eq:14.2.17} must continue to hold if we
interchange \lq\lq in\rq\rq\ and \lq\lq out,\rq\rq\ replace \(C\) and \(D\) by
\(A\) and \(B\), and replace \(S\) by \(S^{-1}\).
\begin{enumerate}[label=\alph*),leftmargin=2em,itemsep=.45ex]
\item Using this modification of equation~\eqref{eq:14.2.17} and the results of
problem~2(b), show that the expected number of particles in state
\(\sigma\in\mathcal H_{\mathrm{out}}\) spontaneously created from the
\lq\lq in\rq\rq\ vacuum state is given by
\[
  \langle N(\sigma)\rangle=\lVert B\sigma\rVert^2.
\]
\item Using equation~\eqref{eq:14.3.6} as well as the result of part~(a) and the
relations proven in problem~3, derive equation~\eqref{eq:14.3.7} for particle
creation by black holes.
\end{enumerate}

\item
The classical Klein--Gordon stress-energy tensor~(4.3.10) has manifestly positive
energy density everywhere; i.e., for any timelike vector \(t^a\) at any point we
always have \(T_{ab}t^at^b\geq0\). Show that this property does \emph{not} carry
over to quantum field theory by finding a state, \(\Psi\), of the free Klein--Gordon
field in Minkowski spacetime for which the normal ordered stress-energy operator
satisfies \(\langle\Psi|\widehat T_{ab}|\Psi\rangle t^at^b<0\) over a region of
spacetime. (Hint: Consider a state \(\Psi\) obtained by superposing the vacuum state
with a small admixture of a two-particle state.) Note, however, that \emph{total}
energy
\[
  E=\int_\Sigma\langle\widehat T_{ab}\rangle\xi^a n^b,
\]
where \(\xi^a\) is a time translation Killing field, always is nonnegative for the free
Klein--Gordon field in Minkowski spacetime.

\item
Consider a box of volume \(V\) filled with blackbody radiation and possibly also
containing a black hole. Ignore the influence of the black hole on the distribution of
radiation in the box. Suppose that the total energy contained in the box (i.e., the mass
of the black hole plus the energy of the radiation) is \(E\).
\begin{enumerate}[label=\alph*),leftmargin=2em,itemsep=.45ex]
\item Write down the formula for the total generalized entropy, \(S'\), of the box as
a function of the mass-energy \(M\) apportioned to the black hole.
\item By extremizing \(S'\) with respect to \(M\), determine the conditions which
must be satisfied by \(E\) and \(V\) in order that equilibrium between the black hole
and radiation be possible (Gibbons and Perry 1978). For \(V=1\text{ meter}^3\),
evaluate the minimum value of \(E\) needed for equilibrium. In the case where
equilibrium is possible, determine which of the extrema of \(S'\) are locally stable by
computing the second derivative of \(S'\).
\item Estimate the conditions on \(E\) and \(V\) such that the configuration of
(globally) maximum generalized entropy will contain a black hole rather than be pure
radiation.
\end{enumerate}
\end{enumerate}
